\section{End-Term University Examination (70 Marks)}

\ExamHeader{End-Term University Examination (70 Marks)}{Paper Code: A (Official Model Paper 1)}{3 Hours}{70 Marks}{-0.25 Marks (Part A)}

\noindent
\textbf{Official University Examination Structure:}
\begin{enumerate}[topsep=2pt,itemsep=1pt]
    \item \textbf{Part A (Objective --- 30 Marks):} 30 MCQs strictly testing \textbf{Unit 4} (Multivariate Diff), \textbf{Unit 5} (Multiple Integrals), and \textbf{Unit 6} (Fourier Series). 10 MCQs per unit, carrying 1 mark each. Negative marking deduction of \textbf{0.25 marks} per incorrect response. All questions in Part A are compulsory.
    \item \textbf{Part B (Descriptive --- 40 Marks):} 5 comprehensive questions carrying \textbf{10 marks each}. Candidates must \textbf{attempt any 4 questions} ($4 \times 10 = 40$ Marks).
    \begin{itemize}
        \item Exactly 3 questions from Units 4, 5, and 6 (one each from Unit 4, Unit 5, Unit 6).
        \item Exactly 2 questions selected from all units (Units 1 to 6).
        \item Sub-parts $(a)\,[5\text{ Marks}] + (b)\,[5\text{ Marks}]$ or single 10-mark problems.
    \end{itemize}
\end{enumerate}
\vspace{3mm}

\subsection*{Part A: Objective Assessment (30 MCQs --- 30 Marks)}

\subsubsection*{Unit 4: Multivariate Differentiation (Questions 1 to 10)}
\mcqitem{1}{If $u = x^3 + y^3 - 3axy$, then $\frac{\partial u}{\partial x}$ is:}{$3x^2 - 3ay$}{$3x^2 - 3ax$}{$3y^2 - 3ax$}{$3x^2$}
\mcqitem{2}{A function $f(x,y)$ is homogeneous of degree $n$ if $f(tx, ty) = $}{$t f(x,y)$}{$t^n f(x,y)$}{$n t f(x,y)$}{$f(x,y)/t^n$}
\mcqitem{3}{By Euler's theorem, if $u$ is homogeneous of degree $n$, then $x u_x + y u_y = $}{$n u$}{$n(n-1)u$}{$u/n$}{$0$}
\mcqitem{4}{The total differential $dz$ of $z = f(x,y)$ is given by:}{$f_x dx + f_y dy$}{$f_x dy + f_y dx$}{$f_{xx} dx + f_{yy} dy$}{$f_x + f_y$}
\mcqitem{5}{At a critical point $(a,b)$, if $D = f_{xx}f_{yy} - (f_{xy})^2 > 0$ and $f_{xx} > 0$, then $(a,b)$ is a:}{Local Minimum}{Local Maximum}{Saddle Point}{Inconclusive}
\mcqitem{6}{At a critical point $(a,b)$, if $D = f_{xx}f_{yy} - (f_{xy})^2 < 0$, then $(a,b)$ is a:}{Local Minimum}{Local Maximum}{Saddle Point}{Inconclusive}
\mcqitem{7}{Clairaut's theorem states that for sufficiently smooth functions:}{$f_{xy} = f_{yx}$}{$f_{xx} = f_{yy}$}{$f_x = f_y$}{$f_{xy} = 0$}
\mcqitem{8}{In the method of Lagrange Multipliers to optimize $f(x,y)$ subject to $g(x,y)=k$, the condition is:}{$\nabla f = \lambda \nabla g$}{$\nabla f \cdot \nabla g = 0$}{$\nabla f = 0$}{$\nabla g = 0$}
\mcqitem{9}{If $z = e^{xy}$, then $\frac{\partial^2 z}{\partial x \partial y}$ is:}{$e^{xy}(1 + xy)$}{$e^{xy}$}{$xy e^{xy}$}{$x^2 y^2 e^{xy}$}
\mcqitem{10}{If $u = \sin(x/y)$, then $x u_x + y u_y$ equals:}{0}{1}{$\sin(x/y)$}{$\cos(x/y)$}

\subsubsection*{Unit 5: Multivariable Integration and Applications (Questions 11 to 20)}
\mcqitem{11}{The Jacobian $J = \frac{\partial(x,y)}{\partial(r,\theta)}$ for polar coordinates $x=r\cos\theta, y=r\sin\theta$ is:}{1}{$r$}{$r^2$}{$1/r$}
\mcqitem{12}{The area of a region $R$ in the $xy$-plane is calculated as:}{$\iint_R 1\,dA$}{$\iint_R xy\,dA$}{$\iint_R (x+y)\,dA$}{$\iint_R (x^2+y^2)\,dA$}
\mcqitem{13}{The volume element $dV$ in cylindrical coordinates $(r, \theta, z)$ is:}{$dr\,d\theta\,dz$}{$r\,dr\,d\theta\,dz$}{$r^2\,dr\,d\theta\,dz$}{$r\sin\theta\,dr\,d\theta\,dz$}
\mcqitem{14}{The volume element $dV$ in spherical polar coordinates $(\rho, \phi, \theta)$ is:}{$\rho^2\sin\phi\,d\rho\,d\phi\,d\theta$}{$\rho\sin\phi\,d\rho\,d\phi\,d\theta$}{$\rho^2\,d\rho\,d\phi\,d\theta$}{$\rho^2\cos\phi\,d\rho\,d\phi\,d\theta$}
\mcqitem{15}{The value of $\int_0^1 \int_0^2 dy\,dx$ is:}{1}{2}{3}{4}
\mcqitem{16}{Changing the order of integration in $\int_0^1 \int_0^x f(x,y)\,dy\,dx$ yields:}{$\int_0^1 \int_y^1 f(x,y)\,dx\,dy$}{$\int_0^1 \int_0^y f(x,y)\,dx\,dy$}{$\int_0^x \int_0^1 f(x,y)\,dx\,dy$}{$\int_0^1 \int_0^1 f(x,y)\,dx\,dy$}
\mcqitem{17}{The value of $\int_0^1 \int_0^1 xy\,dx\,dy$ is:}{1/4}{1/2}{1}{0}
\mcqitem{18}{The solid volume of a sphere of radius $R$ is:}{$\frac{4}{3}\pi R^3$}{$4\pi R^2$}{$\frac{2}{3}\pi R^3$}{$\pi R^3$}
\mcqitem{19}{The integral $\iint_R dx\,dy$ over the circle $x^2+y^2 \le a^2$ represents:}{Circumference}{Planar area $\pi a^2$}{Volume}{Centroid}
\mcqitem{20}{Fubini's theorem states that for continuous $f(x,y)$ on a rectangle:}{Iterated integrals in either order are equal}{Double integral is always zero}{Jacobian is zero}{Order cannot be changed}

\subsubsection*{Unit 6: Introduction to Fourier Series (Questions 21 to 30)}
\mcqitem{21}{The fundamental period of $\cos(nx)$ is:}{$2\pi$}{$2\pi/n$}{$\pi/n$}{$n\pi$}
\mcqitem{22}{For an even function $f(x)$ on $(-\pi, \pi)$, the Fourier coefficient $b_n$ is:}{0}{$a_n$}{$a_0/2$}{$\pi$}
\mcqitem{23}{For an odd function $f(x)$ on $(-\pi, \pi)$, the Fourier coefficients $a_0$ and $a_n$ are:}{Both 0}{Both 1}{$a_0=0, a_n \ne 0$}{Non-zero}
\mcqitem{24}{At a point of jump discontinuity $x_0$, the Fourier series converges to:}{$f(x_0^+)$}{$f(x_0^-)$}{$\frac{f(x_0^+) + f(x_0^-)}{2}$}{0}
\mcqitem{25}{The Euler formula for $a_0$ on $(-\pi, \pi)$ is:}{$\frac{1}{\pi}\int_{-\pi}^\pi f(x)dx$}{$\frac{2}{\pi}\int_{-\pi}^\pi f(x)dx$}{$\frac{1}{2\pi}\int_{-\pi}^\pi f(x)dx$}{$\int_{-\pi}^\pi f(x)dx$}
\mcqitem{26}{The half-range sine series of $f(x)$ on $(0, L)$ contains only:}{Cosine terms}{Sine terms}{Constant term and sines}{Exponential terms}
\mcqitem{27}{The half-range cosine series of $f(x)$ on $(0, L)$ contains only:}{Sine terms}{Cosine terms and constant}{Tangent terms}{Exponential terms}
\mcqitem{28}{Dirichlet conditions require $f(x)$ on $[-L, L]$ to have:}{Infinite discontinuities}{Finite number of maxima, minima and discontinuities}{Continuous derivative everywhere}{Zero integral}
\mcqitem{29}{The Fourier series of $f(x) = x$ on $(-\pi, \pi)$ contains:}{Only cosine terms}{Only sine terms}{Both sine and cosine terms}{No terms}
\mcqitem{30}{Parseval's identity for Fourier series relates the integral of $[f(x)]^2$ to:}{Sum of squares of Fourier coefficients}{Derivative of $f(x)$}{Jacobian determinant}{Trace of matrix}

\newpage
\subsubsection*{Part A Answer Key \& Technical Rationales}

\begin{table}[htbp]
\centering
\caption{\textbf{Part A: Answer Key \& Explanations (Questions 1 to 30)}}
\vspace{2mm}
\begin{tabularx}{\textwidth}{l c X l c X}
\toprule
\textbf{Q\#} & \textbf{Ans} & \textbf{Technical Rationale} & \textbf{Q\#} & \textbf{Ans} & \textbf{Technical Rationale} \\
\midrule
1 & \textbf{(a)} & Partial derivative with respect to $x$ treating $y$ as constant. & 16 & \textbf{(a)} & Reversing slices gives $y \in [0, 1]$ and $x \in [y, 1]$. \\
2 & \textbf{(b)} & Formal definition of degree $n$ homogeneity. & 17 & \textbf{(a)} & $(\int_0^1 x dx)(\int_0^1 y dy) = (1/2)(1/2) = 1/4$. \\
3 & \textbf{(a)} & First-order Euler identity for homogeneous functions. & 18 & \textbf{(a)} & Standard spherical volume formula derived via triple integral. \\
4 & \textbf{(a)} & Definition of total differential $dz$. & 19 & \textbf{(b)} & Double integral of $1$ over circle gives area $\pi a^2$. \\
5 & \textbf{(a)} & Second derivative test criterion for local minimum. & 20 & \textbf{(a)} & Integral order invariance over Cartesian domains. \\
6 & \textbf{(c)} & Negative discriminant strictly defines a saddle point. & 21 & \textbf{(b)} & Period $T = 2\pi / \omega = 2\pi/n$. \\
7 & \textbf{(a)} & Equality of mixed second partial derivatives. & 22 & \textbf{(a)} & Sine terms vanish for even functions. \\
8 & \textbf{(a)} & Collinearity of gradient vectors $\nabla f = \lambda \nabla g$. & 23 & \textbf{(a)} & Cosine terms and constant vanish for odd functions. \\
9 & \textbf{(a)} & $z_x = y e^{xy} \implies z_{xy} = e^{xy} + y(x e^{xy}) = e^{xy}(1+xy)$. & 24 & \textbf{(c)} & Dirichlet convergence criterion at jump discontinuity. \\
10 & \textbf{(a)} & Homogeneous of degree $0 \implies n u = 0 \cdot u = 0$. & 25 & \textbf{(a)} & Standard Euler formula definition. \\
11 & \textbf{(b)} & Determinant of transformation matrix is $r$. & 26 & \textbf{(b)} & Half-range sine series represents odd periodic extension. \\
12 & \textbf{(a)} & Area is the double integral of unit integrand. & 27 & \textbf{(b)} & Half-range cosine series represents even periodic extension. \\
13 & \textbf{(b)} & Cylindrical differential volume is $r\,dr\,d\theta\,dz$. & 28 & \textbf{(b)} & Classical Dirichlet conditions for Fourier convergence. \\
14 & \textbf{(a)} & Spherical Jacobian determinant gives $\rho^2\sin\phi$. & 29 & \textbf{(b)} & $f(x)=x$ is an odd function, hence contains only sine terms. \\
15 & \textbf{(b)} & $(1 - 0)(2 - 0) = 2$. & 30 & \textbf{(a)} & Conservation of energy / harmonic power equation. \\
\bottomrule
\end{tabularx}
\end{table}

\newpage
\subsection*{Part B: Descriptive / Subjective Assessment (40 Marks)}
\noindent
\textbf{Instructions:} Answer any \textbf{FOUR} questions ($4 \times 10 = \mathbf{40\text{ Marks}}$).

\vspace{3mm}
\noindent\textbf{Question 31 (Unit 4: Multivariate Differentiation)} \hfill \textbf{[10 Marks]}
\begin{enumerate}[label=(\alph*),topsep=2pt,itemsep=2pt]
    \item If $u = \sin^{-1}\left(\frac{x+2y+3z}{\sqrt{x^8+y^8+z^8}}\right)$, prove that $x\frac{\partial u}{\partial x} + y\frac{\partial u}{\partial y} + z\frac{\partial u}{\partial z} = -3\tan u$. \hfill \textbf{[5 Marks]}
    \item Find the dimensions of a rectangular box open at the top having maximum volume for a given surface area $S = 108\text{ cm}^2$. \hfill \textbf{[5 Marks]}
\end{enumerate}

\begin{solbox}
\textbf{Part (a):}
Let $w = \sin u = \frac{x+2y+3z}{\sqrt{x^8+y^8+z^8}}$.
Testing homogeneity: $w(tx, ty, tz) = \frac{t(x+2y+3z)}{\sqrt{t^8(x^8+y^8+z^8)}} = \frac{t}{t^4} w(x,y,z) = t^{-3} w(x,y,z)$.
So $w$ is homogeneous of degree $n = -3$.
By Euler's Theorem in 3 variables: $x w_x + y w_y + z w_z = -3 w$.
Since $w = \sin u$, $w_x = \cos u \cdot u_x, w_y = \cos u \cdot u_y, w_z = \cos u \cdot u_z$.
Substituting: $\cos u (x u_x + y u_y + z u_z) = -3\sin u \implies x u_x + y u_y + z u_z = -3\tan u$.

\textbf{Part (b):}
Let dimensions be length $x$, width $y$, height $z$.
Volume $V = xyz$. Surface area of open box $g(x,y,z) = xy + 2yz + 2xz = 108$.
By Lagrange multipliers $\nabla V = \lambda \nabla g$:
\begin{align}
yz &= \lambda(y + 2z) \implies xyz = \lambda(xy + 2xz) \\
xz &= \lambda(x + 2z) \implies xyz = \lambda(xy + 2yz) \\
xy &= \lambda(2y + 2x) \implies xyz = \lambda(2yz + 2xz)
\end{align}
Equating expressions: $xy = 2yz = 2xz \implies x = y = 2z$.
Substitute into constraint: $(2z)(2z) + 2(2z)z + 2(2z)z = 4z^2 + 4z^2 + 4z^2 = 12z^2 = 108 \implies z^2 = 9 \implies z = 3\text{ cm}$.
Therefore, $x = 6\text{ cm}, y = 6\text{ cm}, z = 3\text{ cm}$ with maximum volume $V = 6 \times 6 \times 3 = \mathbf{108\text{ cm}^3}$.
\end{solbox}

\vspace{4mm}
\noindent\textbf{Question 32 (Unit 5: Multiple Integrals)} \hfill \textbf{[10 Marks]}
\begin{enumerate}[label=(\alph*),topsep=2pt,itemsep=2pt]
    \item Change the order of integration and evaluate $I = \int_0^1 \int_x^{\sqrt{2-x^2}} \frac{x}{\sqrt{x^2+y^2}}\,dy\,dx$. \hfill \textbf{[5 Marks]}
    \item Evaluate the triple integral $\iiint_V z\,dx\,dy\,dz$ over the region $V$ bounded by the surfaces $z = 0$, $z = x^2+y^2$, and $x^2+y^2 = 4$. \hfill \textbf{[5 Marks]}
\end{enumerate}

\begin{solbox}
\textbf{Part (a):}
In polar coordinates: $x = r\cos\theta, y = r\sin\theta$.
The bounding curves are $y = x \implies \theta = \pi/4$, and $y = \sqrt{2-x^2} \implies x^2+y^2 = 2 \implies r = \sqrt{2}$.
As $x$ goes from $0$ to $1$, the region covers $\theta \in [\pi/4, \pi/2]$ with $r \in [0, \sqrt{2}]$.
\begin{align}
I &= \int_{\pi/4}^{\pi/2} \int_0^{\sqrt{2}} \frac{r\cos\theta}{r} r\,dr\,d\theta = \left( \int_{\pi/4}^{\pi/2} \cos\theta\,d\theta \right) \left( \int_0^{\sqrt{2}} r\,dr \right) \\
&= \Big[ \sin\theta \Big]_{\pi/4}^{\pi/2} \left[ \frac{r^2}{2} \right]_0^{\sqrt{2}} = \left( 1 - \frac{1}{\sqrt{2}} \right) (1) = \mathbf{1 - \frac{1}{\sqrt{2}} = \frac{\sqrt{2}-1}{\sqrt{2}}}
\end{align}

\textbf{Part (b):}
Transform to cylindrical coordinates: $x = r\cos\theta, y = r\sin\theta, z = z$.
$r \in [0, 2]$, $\theta \in [0, 2\pi]$, $z \in [0, r^2]$.
\begin{align}
\iiint_V z\,dV &= \int_0^{2\pi} \int_0^2 \int_0^{r^2} z \cdot r\,dz\,dr\,d\theta = (2\pi) \int_0^2 r \left[ \frac{z^2}{2} \right]_0^{r^2} dr \\
&= (2\pi) \int_0^2 \frac{r^5}{2}\,dr = \pi \left[ \frac{r^6}{6} \right]_0^2 = \pi \left( \frac{64}{6} \right) = \mathbf{\frac{32\pi}{3}}
\end{align}
\end{solbox}

\vspace{4mm}
\noindent\textbf{Question 33 (Unit 6: Fourier Series)} \hfill \textbf{[10 Marks]}
\begin{enumerate}[label=(\alph*),topsep=2pt,itemsep=2pt]
    \item Obtain the Fourier series expansion of $f(x) = x - x^2$ in the interval $(-\pi, \pi)$. \hfill \textbf{[5 Marks]}
    \item Find the half-range cosine series for $f(x) = (x-1)^2$ in the interval $0 < x < 1$. \hfill \textbf{[5 Marks]}
\end{enumerate}

\begin{solbox}
\textbf{Part (a):}
Decompose $f(x) = x - x^2$ into odd part $f_{\text{odd}}(x) = x$ and even part $f_{\text{even}}(x) = -x^2$.
\begin{itemize}
    \item $a_0 = \frac{1}{\pi}\int_{-\pi}^\pi -x^2 dx = -\frac{2\pi^2}{3}$.
    \item $a_n = \frac{1}{\pi}\int_{-\pi}^\pi -x^2 \cos(nx)dx = \frac{4(-1)^{n+1}}{n^2}$.
    \item $b_n = \frac{1}{\pi}\int_{-\pi}^\pi x \sin(nx)dx = \frac{2(-1)^{n+1}}{n}$.
\end{itemize}
Thus:
\begin{equation}
f(x) = -\frac{\pi^2}{3} + \sum_{n=1}^\infty \left[ \frac{4(-1)^{n+1}}{n^2}\cos(nx) + \frac{2(-1)^{n+1}}{n}\sin(nx) \right]
\end{equation}

\textbf{Part (b):}
Here $L = 1$.
$a_0 = \frac{2}{1}\int_0^1 (x-1)^2 dx = 2\left[ \frac{(x-1)^3}{3} \right]_0^1 = 2(0 - (-1/3)) = \frac{2}{3}$.
\begin{align}
a_n &= 2\int_0^1 (x-1)^2 \cos(n\pi x)dx = 2 \left[ (x-1)^2 \frac{\sin n\pi x}{n\pi} - 2(x-1)\frac{-\cos n\pi x}{n^2\pi^2} + 2\frac{-\sin n\pi x}{n^3\pi^3} \right]_0^1 \\
&= 2 \left[ 0 - \frac{2}{n^2\pi^2} \right] = \frac{4}{n^2\pi^2}
\end{align}
Therefore, the half-range cosine series is:
\begin{equation}
f(x) = \frac{1}{3} + \frac{4}{\pi^2}\sum_{n=1}^\infty \frac{1}{n^2}\cos(n\pi x)
\end{equation}
\end{solbox}

\vspace{4mm}
\noindent\textbf{Question 34 (Unit 1 \& Unit 2: Linear Algebra \& Calculus)} \hfill \textbf{[10 Marks]}
\begin{enumerate}[label=(\alph*),topsep=2pt,itemsep=2pt]
    \item Using the Cayley--Hamilton Theorem, find the inverse of $A = \begin{bmatrix} 1 & 1 & 3 \\ 1 & 3 & -3 \\ -2 & -4 & -4 \end{bmatrix}$. \hfill \textbf{[5 Marks]}
    \item If $y = (x^2-1)^n$, prove that $(x^2-1)y_{n+2} + 2xy_{n+1} - n(n+1)y_n = 0$. \hfill \textbf{[5 Marks]}
\end{enumerate}

\begin{solbox}
\textbf{Part (a):}
$S_1 = 1 + 3 - 4 = 0$.
$S_2 = (-12-12) + (-4+6) + (3-1) = -24 + 2 + 2 = -20$.
$S_3 = \det(A) = 1(-24) - 1(-10) + 3(2) = -24 + 10 + 6 = -8$.
Characteristic equation: $\lambda^3 - 20\lambda + 8 = 0$.
By Cayley--Hamilton: $A^3 - 20A + 8I = O \implies A^{-1} = -\frac{1}{8}(A^2 - 20I)$.
Computing $A^2$:
\begin{equation}
A^2 = \begin{bmatrix} -4 & -8 & -12 \\ 10 & 22 & 6 \\ 2 & 2 & 22 \end{bmatrix} \implies A^2 - 20I = \begin{bmatrix} -24 & -8 & -12 \\ 10 & 2 & 6 \\ 2 & 2 & 2 \end{bmatrix}
\end{equation}
Thus, $A^{-1} = \frac{1}{8}\begin{bmatrix} 24 & 8 & 12 \\ -10 & -2 & -6 \\ -2 & -2 & -2 \end{bmatrix} = \begin{bmatrix} 3 & 1 & 3/2 \\ -5/4 & -1/4 & -3/4 \\ -1/4 & -1/4 & -1/4 \end{bmatrix}$.

\textbf{Part (b):}
$y = (x^2-1)^n \implies y_1 = 2nx(x^2-1)^{n-1} \implies (x^2-1)y_1 = 2nxy$.
Differentiating $(n+1)$ times by Leibniz theorem:
\begin{align}
y_{n+2}(x^2-1) + (n+1)y_{n+1}(2x) + \frac{(n+1)n}{2}y_n(2) &= 2n[y_{n+1}x + (n+1)y_n] \\
(x^2-1)y_{n+2} + 2xy_{n+1} - n(n+1)y_n &= 0
\end{align}
\end{solbox}

\vspace{4mm}
\noindent\textbf{Question 35 (Unit 3 \& Unit 5: Integrals)} \hfill \textbf{[10 Marks]}
\begin{enumerate}[label=(\alph*),topsep=2pt,itemsep=2pt]
    \item Evaluate $\int_0^{\pi/2} \ln(\sin x)\,dx$. \hfill \textbf{[5 Marks]}
    \item Find the area of the region enclosed between the parabolas $y^2 = 4ax$ and $x^2 = 4ay$ using double integration. \hfill \textbf{[5 Marks]}
\end{enumerate}

\begin{solbox}
\textbf{Part (a):}
Let $I = \int_0^{\pi/2} \ln(\sin x)\,dx = \int_0^{\pi/2} \ln(\cos x)\,dx$.
\begin{align}
2I &= \int_0^{\pi/2} \ln(\sin x \cos x)\,dx = \int_0^{\pi/2} \ln\left(\frac{\sin 2x}{2}\right)\,dx \\
&= \int_0^{\pi/2} \ln(\sin 2x)\,dx - \ln 2 \int_0^{\pi/2} 1\,dx = \frac{1}{2}\int_0^\pi \ln(\sin t)\,dt - \frac{\pi}{2}\ln 2 \\
&= \int_0^{\pi/2} \ln(\sin t)\,dt - \frac{\pi}{2}\ln 2 = I - \frac{\pi}{2}\ln 2 \implies I = \mathbf{-\frac{\pi}{2}\ln 2}
\end{align}

\textbf{Part (b):}
Intersection points: $\left(\frac{x^2}{4a}\right)^2 = 4ax \implies x^4 = 64a^3 x \implies x = 0, \; x = 4a$.
For $x \in [0, 4a]$, $y$ ranges from $\frac{x^2}{4a}$ up to $2\sqrt{a}\sqrt{x}$.
\begin{align}
\text{Area } A &= \int_0^{4a} \int_{x^2/(4a)}^{2\sqrt{a}\sqrt{x}} dy\,dx = \int_0^{4a} \left( 2\sqrt{a}x^{1/2} - \frac{x^2}{4a} \right) dx \\
&= \left[ 2\sqrt{a} \frac{2x^{3/2}}{3} - \frac{x^3}{12a} \right]_0^{4a} = \frac{4\sqrt{a}}{3}(8a^{3/2}) - \frac{64a^3}{12a} = \frac{32a^2}{3} - \frac{16a^2}{3} = \mathbf{\frac{16a^2}{3}}
\end{align}
\end{solbox}
